\documentclass[11pt,a4paper]{article}
\usepackage[utf8]{inputenc}
\usepackage[T1]{fontenc}
\usepackage{lmodern}
\usepackage{graphicx}
\usepackage{svg}

% Load actor data commands
% Provenance: Six-Axis Political Compass v2.6.1 (718eeae82bd1affecdcba27103e5bc453b327cfe) generated 2026-05-20T13:05:52.110Z
% Source: data/actors/Conservative-Party.json
\newcommand{\ConservativePartyColor}{\#1A75BB}
\newcommand{\ConservativePartyCultural}{7.0}
\newcommand{\ConservativePartyEconomic}{3.0}
\newcommand{\ConservativePartyMilitary}{8.0}
\newcommand{\ConservativePartySovereignty}{6.0}
\newcommand{\ConservativePartyGovernance}{5.0}
\newcommand{\ConservativePartyClass}{1.0}

% Provenance: Six-Axis Political Compass v2.2.3 (76548cba7849e851938b9bd62b6adf23acbdf64c) generated 2026-05-19T10:45:46.188Z
% Source: data/actors/Labour-Party.json
\newcommand{\LabourPartyColor}{\#c0392b}
\newcommand{\LabourPartyCultural}{5.0}
\newcommand{\LabourPartyEconomic}{3.0}
\newcommand{\LabourPartyMilitary}{7.0}
\newcommand{\LabourPartySovereignty}{5.0}
\newcommand{\LabourPartyGovernance}{6.0}
\newcommand{\LabourPartyClass}{3.0}

% Provenance: Six-Axis Political Compass v1.0.0 (6d0121db747528bec94997d0427e57823672befd) generated 2026-05-16T13:15:49.307Z
% Source: src/data.js — actor "Winston Churchill"
\newcommand{\WinstonChurchillColor}{\#1A237E}
\newcommand{\WinstonChurchillCultural}{8.0}
\newcommand{\WinstonChurchillSovereignty}{8.0}
\newcommand{\WinstonChurchillMilitary}{9.0}
\newcommand{\WinstonChurchillEconomic}{3.0}
\newcommand{\WinstonChurchillClass}{2.0}
\newcommand{\WinstonChurchillLiberty}{3.0}


\title{The Six-Axis Political Compass}
\author{A Common Enemy Project}
\date{\today}

\begin{document}

\maketitle

\section{Introduction}

The Six-Axis Political Compass maps political actors across six dimensions:
Cultural, Economic, Military, Sovereignty, Governance, and Class.

\section{Individual Profiles}

The Conservative Party scores
\ConservativePartyCultural{} on the Cultural axis,
\ConservativePartyEconomic{} on the Economic axis, and
\ConservativePartyMilitary{} on the Military axis.

\begin{figure}[h]
  \centering
  \includesvg[width=0.45\textwidth]{actors/Conservative-Party.svg}
  \caption{Radar chart of the Conservative Party.}
  \label{fig:conservative}
\end{figure}

Winston Churchill scores
\WinstonChurchillCultural{} on Cultural and
\WinstonChurchillMilitary{} on Military.

\begin{figure}[h]
  \centering
  \includesvg[width=0.45\textwidth]{actors/Winston-Churchill.svg}
  \caption{Radar chart of Winston Churchill.}
  \label{fig:churchill}
\end{figure}

\section{Comparison Charts}

Figure~\ref{fig:uk-parties} overlays the major UK parties.

\begin{figure}[h]
  \centering
  \includesvg[width=0.7\textwidth]{comparisons/uk-parties.svg}
  \caption{UK Political Parties on the Six-Axis Compass.}
  \label{fig:uk-parties}
\end{figure}

\end{document}
